\documentclass{article}
\usepackage{amsmath, amssymb, graphicx, natbib}

\begin{document}

\title{Modelo Matemático para la Predicción de la Generación de CO$_2$ en la Bioconversión de Residuos por Larvas de Mosca Soldado Negra}
\author{}
\date{\today}
\maketitle

\section{Introducción}
El proceso de bioconversión de residuos orgánicos mediante larvas de la mosca soldado negra (\textit{Hermetia illucens}, BSFL) genera CO$_2$ como subproducto de la respiración larvaria, la actividad microbiana y la biodegradación aeróbica y anaeróbica de la materia orgánica. Este modelo matemático permite predecir la cantidad de CO$_2$ generada en función de la cantidad inicial de huevos depositados sobre un sustrato, considerando la dinámica de crecimiento larvario y las interacciones ambientales, especialmente la humedad y la temperatura.

\section{Variables y Parámetros del Modelo}
\subsection{Variables dependientes del tiempo $t$}
\begin{itemize}
    \item $N(t)$: Número de larvas en el sistema.
    \item $X(t)$ [g]: Biomasa larval.
    \item $S(t)$ [g]: Biomasa de residuos disponibles.
    \item $C(t)$ [g]: Cantidad acumulada de CO$_2$ producido.
    \item $H_s(t)$ [\%]: Humedad de la biomasa.
    \item $H_a(t)$ [\%]: Humedad del ambiente.
    \item $O_2(t)$ [\%]: Concentración de oxígeno en el sistema.
\end{itemize}

\subsection{Parámetros constantes}
\begin{itemize}
    \item $\mu_{\max}$ [1/día]: Tasa máxima de crecimiento larval.
    \item $K_s$ [g]: Constante de semi-saturación del sustrato (Monod).
    \item $Y_{X/S}$ [g larva/g residuo]: Rendimiento de conversión de sustrato en biomasa larvaria.
    \item $Y_{C/S}$ [g CO$_2$/g residuo]: Rendimiento de CO$_2$ por unidad de residuo consumido.
    \item $\lambda_H$ [1/día]: Coeficiente de transferencia de humedad entre la biomasa y el ambiente.
    \item $k_O$ [1/día]: Tasa de consumo de oxígeno por unidad de biomasa larvaria.
    \item $E_H$ [1/día]: Tasa de evaporación de humedad desde la biomasa.
\end{itemize}

\section{Ecuaciones del Modelo}
\subsection{Dinámica de Población Larval}
\begin{equation}
    N(t) = \eta N_0 e^{-\delta t}
\end{equation}

donde $\eta$ es la tasa de eclosión y $\delta$ es la tasa de mortalidad larval.

\subsection{Crecimiento Larvario}
\begin{equation}
    \frac{dX}{dt} = Y_{X/S} \frac{\mu_{\max} S}{K_s + S} X \left( 1 - \frac{X}{X_{\max}} \right) - \delta X
\end{equation}

donde $X_{\max}$ es la biomasa larvaria máxima sostenible.

\subsection{Consumo de Residuos}
\begin{equation}
    \frac{dS}{dt} = -\frac{\mu_{\max} S}{K_s + S} X - k_s S
\end{equation}

donde $k_s$ representa la degradación espontánea del sustrato no consumido.

\subsection{Producción de CO$_2$}
\begin{equation}
    \frac{dC}{dt} = Y_{C/S} \frac{\mu_{\max} S}{K_s + S} X + k_M S \frac{O_2}{K_O + O_2}
\end{equation}

donde $k_M$ es la tasa de producción de CO$_2$ por metabolismo microbiano y $K_O$ es la constante de semi-saturación de oxígeno.

\subsection{Dinámica de Humedad y Oxígeno}
\begin{equation}
    \frac{dH_s}{dt} = -\lambda_H (H_s - H_a) - E_H H_s
\end{equation}

\begin{equation}
    \frac{dH_a}{dt} = \lambda_H (H_s - H_a)
\end{equation}

\begin{equation}
    \frac{dO_2}{dt} = -k_O X - k_M S \frac{O_2}{K_O + O_2}
\end{equation}

donde el oxígeno es consumido tanto por las larvas como por los microorganismos presentes en la biodegradación aeróbica.

\section{Interpretación del Modelo y Aplicaciones}
Este modelo matemático permite:
\begin{itemize}
    \item Predecir la producción de CO$_2$ en función de la cantidad inicial de huevos y la biomasa de sustrato.
    \item Analizar la influencia de la humedad en la eficiencia del proceso.
    \item Optimizar el diseño de biorreactores y sistemas de cultivo, ajustando condiciones ambientales para maximizar la conversión de residuos y minimizar la emisión de gases de efecto invernadero.
\end{itemize}

La implementación de este modelo en sistemas de monitoreo IoT permitirá obtener datos en tiempo real y ajustar parámetros ambientales automáticamente, favoreciendo un enfoque de economía circular y mitigación del cambio climático.

\bibliographystyle{apalike}
\bibliography{referencias}

\end{document}
